\chapter*{Abstract}
\thispagestyle{empty}

\section*{English Abstract}

This project presents the design, implementation, and deployment of an Optical Character
Recognition (OCR) system for printed text using deep learning techniques. The system
employs a Convolutional Recurrent Neural Network (CRNN) architecture combining
Convolutional Neural Networks (CNN) for spatial feature extraction, Bidirectional
Long Short-Term Memory (BiLSTM) networks for temporal sequence modeling, and
Connectionist Temporal Classification (CTC) loss for end-to-end training without
explicit character alignment.

The dataset comprises 422,400 grayscale character images across 66 classes (digits,
uppercase and lowercase letters, and punctuation marks), with 6,400 images per class.
A comprehensive preprocessing pipeline including Otsu binarization, aspect-ratio-preserving
resizing, and normalization ensures robust input quality.

The trained model is deployed as a Flask web application providing real-time character
recognition through a user-friendly interface. The system implements automatic character
segmentation via contour detection for multi-character text recognition.

Key technical contributions include handling legacy HDF5 model formats in modern Keras 3
environments, implementing efficient data pipelines with tf.data, and deploying a
production-ready web interface with CORS support.

Evaluation reveals functional character recognition with identified areas for improvement,
including architecture optimization (switching from sequence models to classification
for single-character recognition) and data augmentation for enhanced robustness.

\keywordss{Optical Character Recognition, Deep Learning, Convolutional Neural Networks,
Recurrent Neural Networks, LSTM, CTC Loss, TensorFlow, Keras, Flask, Computer Vision}

\vspace{2cm}

\section*{Résumé en Français}

Ce projet présente la conception, l'implémentation et le déploiement d'un système de
Reconnaissance Optique de Caractères (OCR) pour texte imprimé utilisant des techniques
d'apprentissage profond. Le système emploie une architecture de Réseau de Neurones
Convolutif Récurrent (CRNN) combinant des Réseaux de Neurones Convolutifs (CNN) pour
l'extraction de caractéristiques spatiales, des réseaux Long Short-Term Memory
Bidirectionnels (BiLSTM) pour la modélisation de séquences temporelles, et la perte
Connectionist Temporal Classification (CTC) pour l'entraînement de bout en bout sans
alignement explicite des caractères.

Le jeu de données comprend 422 400 images de caractères en niveaux de gris réparties
sur 66 classes (chiffres, lettres majuscules et minuscules, et signes de ponctuation),
avec 6 400 images par classe. Un pipeline de prétraitement complet incluant la
binarisation d'Otsu, le redimensionnement préservant le ratio d'aspect et la
normalisation assure une qualité d'entrée robuste.

Le modèle entraîné est déployé comme application web Flask fournissant une reconnaissance
de caractères en temps réel via une interface conviviale. Le système implémente une
segmentation automatique des caractères via détection de contours pour la reconnaissance
de texte multi-caractères.

Les contributions techniques clés incluent la gestion des formats de modèles HDF5
hérités dans les environnements Keras 3 modernes, l'implémentation de pipelines de
données efficaces avec tf.data, et le déploiement d'une interface web prête pour la
production avec support CORS.

L'évaluation révèle une reconnaissance fonctionnelle des caractères avec des domaines
d'amélioration identifiés, notamment l'optimisation de l'architecture (passage des
modèles de séquence à la classification pour la reconnaissance de caractères uniques)
et l'augmentation de données pour une robustesse accrue.

\keywords{Reconnaissance Optique de Caractères, Apprentissage Profond, Réseaux de
Neurones Convolutifs, Réseaux de Neurones Récurrents, LSTM, Perte CTC, TensorFlow,
Keras, Flask, Vision par Ordinateur}

\end{document}
